\section{The finite parity calculus}\label{sec:parity}

\subsection{A definition that requires only rational comparisons}

Let $P=(v_0,\ldots,v_{s-1},v_s=v_0)$ be a rational polygonal
closed walk. Repetitions, backtracking, and self-intersections are allowed.
Let $p\in\Q^2$ be off every edge. For the edge $[a,b]$, put
\[
 \sigma_p(a,b)=[a_y>p_y]\xor[b_y>p_y].
\]
Here brackets turn a decidable statement into a bit. If $\sigma_p=0$,
give the edge contribution zero. If $\sigma_p=1$, the denominator
$b_y-a_y$ is nonzero and the horizontal intersection coordinate is
\[
 x_* = a_x+(p_y-a_y)\frac{b_x-a_x}{b_y-a_y}.
\]
Give the edge contribution $c_p(a,b)=[x_*>p_x]$. Define
\begin{equation}
 \parity(P,p)=\bigoplus_{i<s}c_p(v_i,v_{i+1})\in\F.
 \label{eq:parity}
\end{equation}
Division can be removed by multiplying the comparison by the denominator
and inspecting its sign. The half-open height convention deals with a
ray passing exactly through a vertex. No test of equality of real
numbers is present.

The following facts are immediate from the definition: contributions
are unchanged by reversing an edge; the xor over concatenated edge
lists is the xor of their values; and two copies of the same edge cancel.
In particular, we may work with finite edge chains modulo two. A chain
may consist of several closed walks, and its parity is their xor.

\begin{lemma}[The triangle calculation]\label{lem:triangle}
If $p$ is outside the closed convex hull of a rational triangle, the
parity of its boundary about $p$ is zero. If $Q$ is a nondegenerate
convex rational square or parallelogram, its parity is one at any
interior point and zero at any point outside its closed region.
\end{lemma}

\begin{proof}
Intersect the horizontal line through $p$ with the triangle.
Using the three linear edge equations, this intersection is empty,
a point, or an interval. If $p$ is outside the triangle, all the
intersection coordinates lie on the same side of $p_x$, or there
are no intersection coordinates. If they are all left, every
contribution is zero. If all are right, the xor of the contributions is
\[
 ([v_{0y}>p_y]\xor[v_{1y}>p_y])
 \xor([v_{1y}>p_y]\xor[v_{2y}>p_y])
 \xor([v_{2y}>p_y]\xor[v_{0y}>p_y])=0.
\]
At a lowest vertex the two incident straddles are both counted; at a
highest vertex neither is counted; at a middle-height vertex exactly
one is counted. Thus the same calculation covers height ties and
horizontal edges. Degenerate triangles are covered too: the horizontal
slice is still an interval or a point, and the same xor cancels.

For a convex square or parallelogram and an interior point, the
horizontal slice has one endpoint strictly to the left and one strictly
to the right. At each endpoint the half-open rule counts one incident
edge modulo two. Precisely the right endpoint contributes to the ray,
so the answer is one. At an exterior point the preceding all-left or
all-right argument gives zero. These assertions can equivalently be
checked by ordering the finitely many vertex heights and solving
linear inequalities; they invoke no theorem about arbitrary polygons.
\end{proof}

\begin{lemma}[Finite replacement]\label{lem:replacement}
Suppose two rational closed edge chains differ, modulo two, by the
boundaries of finitely many triangles whose closed convex hulls avoid
$p$. Their parities about $p$ are equal.
\end{lemma}

\begin{proof}
Take the xor of the triangle identities in Lemma~\ref{lem:triangle}.
Every internal edge occurs twice and cancels. What remains is the
two stated boundary chains. This is a finite sum identity.
\end{proof}

An important special case replaces an edge $[u,v]$ by $[u,w],[w,v]$
when the triangle $uvw$ avoids the query. Collinear subdivision is
another special case.

\begin{lemma}[Moving the query]\label{lem:moving}
If the segment $[p,q]$ avoids every edge of a rational closed walk $P$,
then $\parity(P,p)=\parity(P,q)$. The same holds when $p,q$ are joined
by a finite polygonal path avoiding $P$.
\end{lemma}

\begin{proof}
Translation changes neither the rational crossing test nor its parity:
$\parity(P,p)=\parity(P-p,0)$. For each edge $[u,v]$ of $P$, compare
the edges $[u-p,v-p]$ and $[u-q,v-q]$. The quadrilateral between them
consists of the points
\[
           z-w,\qquad z\in[u,v],\quad w\in[p,q].
\]
It misses zero, since the two segments do not meet. Its two
triangles therefore miss zero. Applying
Lemma~\ref{lem:replacement} to all these quadrilaterals cancels the
connecting edges at consecutive vertices and proves the equality.
For a polygonal query path, repeat the equality for its successive
segments.
\end{proof}

For rational finite segments, disjointness in this lemma has a rational
positive-distance certificate: minimize the relevant piecewise-linear
distance by finite rational arithmetic. In its later applications,
the certificate also follows directly from the displayed geometric
margins.

\subsection{Stable parity of a named closed curve}

\begin{proposition}[Sampling and local constancy]\label{prop:stable}
Let a named closed curve have a forward modulus, and let $x$ have
clearance at least $r=2^{-k}$. With
\[
 \ell=k+5,\quad
 P_\ell=(G(j/N(\ell),\ell))_{j<N(\ell)},\quad
 x_\ell=X(\ell),
\]
define
\begin{equation}
 b_\Gamma(x)=\parity(P_\ell,x_\ell).
 \label{eq:curvebit}
\end{equation}
This definition is independent of the clearance certificate, sufficiently
fine meshes, and approximation choices. If
$\norm{x-y}<r/4$, then $y$ is off the curve and
$b_\Gamma(y)=b_\Gamma(x)$. More generally, a uniformly continuous
named query path with a supplied positive clearance from the curve
has equal endpoint bits.
\end{proposition}

\begin{proof}
On a parameter mesh interval the curve moves by at most $r/32$,
and each vertex approximation costs at most $r/32$.
Each polygon edge is therefore in the maximum-norm ball of radius
$r/16$ around the curve value at the interval's left endpoint.
The query approximation has error at most $r/32$, so it is off all
the polygon edges.

First compare two meshes both satisfying these error bounds.
Interleave their rational parameter lists. Inserting an intermediate
vertex changes the closed edge list by a triangle. All its vertices
are in a ball of radius at most $r/8$ about one curve value,
after allowing for the errors in both lists. This ball avoids the
query approximation. Matching two approximations at a common
parameter also gives triangles in the same ball. Finite replacement
proves that the parities coincide.

If a coarser valid certificate uses a different $r$, compare both
calculations with one mesh fine enough for the smaller margin.
The same refinement argument applies to each. Query approximations
are compared by a short query segment, which is separated from the
sampled curve; use Lemma~\ref{lem:moving}.

When $\norm{x-y}<r/4$, the segment from $x$ to $y$ stays at least
$3r/4$ from every curve point. Approximate it and the curve at a
precision less than $r/32$ and use Lemma~\ref{lem:moving}. This
also proves the explicit local constancy assertion.

Finally, suppose a path $\alpha$ has clearance $\rho>0$. Its supplied
modulus gives a finite parameter subdivision on which the image
diameter is below $\rho/16$. Approximate its vertices, and the curve,
within $\rho/16$. The polygonal path so obtained stays positively
away from the sampled closed curve. Apply the finite lemma along its
edges, then the already proved independence of approximations at
its endpoints. Every subdivision length was chosen from the
specified modulus and clearance.
\end{proof}

The proposition applies with any fixed rational closed walk in place
of $\Gamma$, and a named path moving the query. This observation will
be used for curve half-arcs and for a rational square. It only needs
forward continuity and clearance; injectivity is not used.

\subsection{A short-loop factorization lemma}
\label{sec:factor}

The inverse modulus enters at exactly this point. It turns geometric
closeness into closeness on the parameter circle.

\begin{lemma}[Parity factors through the parameter circle]\label{lem:factor}
Fix $\delta=2^{-k}$, and put
\begin{equation}
 r=U(k+4),\qquad
 \eta=2^{-V(r)},\qquad H=V(r)+10,\qquad h=2^{-H}=\eta/1024.
 \label{eq:scale}
\end{equation}
Let $P=(x_0,\ldots,x_{s-1},x_s=x_0)$ be a rational closed walk
whose edges have maximum-norm length at most $2h$, and whose
vertices are each within $16h$ of $\Gamma$, with the near points
available by finite sampling. There is one bit $m_P$ such that
\begin{equation}
 \parity(P,a)=m_P\,b_\Gamma(a)
 \label{eq:factor}
\end{equation}
for every $a$ of clearance at least $2\delta$ from $\Gamma$.
The bit is independent of $a$. Simplicity of $P$ is not required.
\end{lemma}

\begin{proof}
Use the uniform cyclic mesh with $N=N(H+4)$ and
$Y_j=G(j/N,H+4)$. Its net error is at most $h/8$ and its evaluation
error is at most $h/16$. For each $x_i$, choose the nearest $Y_j$,
breaking ties by the least index, and call that index $j_i$.
The assumed nearby curve point and the net estimate show
\[
 \norm{x_i-\gamma(j_i/N)}<20h.
\]
This deliberately generous inequality leaves room for all
approximations. Hence
\begin{equation}
 \norm{\gamma(j_i/N)-\gamma(j_{i+1}/N)}
        <20h+2h+20h=42h<\eta.
 \label{eq:42h}
\end{equation}
By the inverse condition and the decidability of the rational parameter
comparison,
\[
 d_\T(j_i/N,j_{i+1}/N)<2^{-r}<1/4.
\]
There is consequently a unique shorter arc between these mesh
parameters, obtainable by comparing their integer indices modulo $N$.
Every intermediate parameter on this arc is at circular distance
less than $2^{-r}$ from its starting parameter. By the forward modulus,
its curve value is within $\delta/16$ of the starting curve value.

Replace $[x_i,x_{i+1}]$ first by the chain through
$Y_{j_i},Y_{j_{i+1}}$, then by the consecutive mesh edges on the
short parameter arc. All replacement triangles lie within
$\delta/8$ of $\gamma(j_i/N)$: their vertices cost at most
$20h$, $\delta/16$, and $h/16$, and
\[
             20h+\delta/16+h/16<\delta/8
\]
by the normalization of $U,V$ in \eqref{eq:scale}.
They avoid every query with clearance $2\delta$.
Finite replacement therefore preserves parity simultaneously for
all the queries in the statement.

After the replacements, we have a closed walk $W$ on the finite
cycle with vertices $0,\ldots,N-1$. Count each undirected cycle edge
modulo two, and write its usage bit as $e_j$. At vertex $j$,
the number of occurrences of entering or leaving is even, because
the walk is closed. Thus
\[
                  e_{j-1}\xor e_j=0.
\]
Following the cyclic indices gives $e_0=e_1=\cdots=e_{N-1}=m_P$.
The parity sum for $W$ is accordingly $m_P$ times the parity sum
for the full sample cycle. The latter is $b_\Gamma(a)$ by
Proposition~\ref{prop:stable}, proving \eqref{eq:factor}.

All counts are finite. Each replaced short arc has at most $N$
edges, so the expanded walk has at most $sN$ edges. The edge-usage
bits depend on $P$ and its selected mesh labels, not on the query.
\end{proof}

This is the decisive arithmetic mechanism: a loop following the curve
closely cannot have arbitrary unrelated parities at distant points.
Its parity pattern is either zero everywhere under consideration or
the curve's own pattern.
